\section{Dimension and depth of Noetherian local rings}
\label{section:0.16}

\subsection{Dimension of a ring}
\label{subsection:0.16.1}
\label{0.16.1}

\begin{env}[16.1.1]
\label{0.16.1.1}
We call the \emph{dimension} (or \emph{Krull dimension}) of a ring $A$ and denote by $\dim A$ the (combinatorial) dimension of its spectrum $\Spec(A)$ \sref{0.14.2.1}; as the irreducible closed subsets of $\Spec(A)$ are the $V(\mathfrak{p})$, where $\mathfrak{p}$ is a prime ideal of $A$ (Bourbaki, \emph{Alg.~comm.}, chap.~II, \textsection~4, n\textsuperscript{o}~3, cor.~1 of prop.~14), it amounts to saying that $\dim(A)$ is \emph{the upper bound of the lengths of chains} \sref{0.14.1.1}
\[
  \mathfrak{p}_0\subset\mathfrak{p}_1\subset\ldots\subset\mathfrak{p}_n
\]
of \emph{prime ideals} of $A$.
If $A\neq 0$, we have $\dim A\geq 0$.
A nonzero Artinian ring (in particular a field) is of dimension $0$ (Bourbaki, \emph{Alg.~comm.}, chap.~IV, \textsection~2, n\textsuperscript{o}~5, prop.~7); a Dedekind ring that is not a field (in particular $\mathbf{Z}$ or a polynomial ring $k[T]$ over a field $k$) is of dimension 1 (Bourbaki, \emph{Alg.~comm.}, chap.~VII, \textsection~2).
A Noetherian ring is not necessarily of finite dimension [30].
If $(\mathfrak{p}_\alpha)$ is the family of minimal prime ideals of $A$, we have \sref{0.14.1.2.1}
\begin{equation}
\label{0.16.1.1.1}
\tag{16.1.1.1}
  \dim(A)=\sup_\alpha(\dim(A/\mathfrak{p}_\alpha)).
\end{equation}
\end{env}

\begin{env}[16.1.2]
\label{0.16.1.2}
For every ideal $\mathfrak{a}$ of $A$, the prime ideals of $A/\mathfrak{a}$ correspond bijectively to the prime ideals of $A$ containing $\mathfrak{a}$, hence
\begin{equation}
\label{0.16.1.2.1}
\tag{16.1.2.1}
  \dim(A/\mathfrak{a})\leq\dim(A).
\end{equation}
If in addition $\mathfrak{a}$ \emph{is not contained in any minimal prime ideal} $\mathfrak{p}_\alpha$ of $A$, every chain of prime ideals of $A$ containing $\mathfrak{a}$ can be completed by one of the $\mathfrak{p}_\alpha$, hence
\begin{equation}
\label{0.16.1.2.2}
\tag{16.1.2.2}
  1+\dim(A/\mathfrak{a})\leq\dim(A).
\end{equation}
\end{env}

\begin{env}[16.1.3]
\label{0.16.1.3}
\oldpage[0\textsubscript{IV-1}]{23}
If $S$ is a multiplicative subset of $A$, the prime ideals of $S^{-1}A$ correspond bijectively to the prime ideals of $A$ not meeting $S$, hence
\begin{equation}
\label{0.16.1.3.1}
\tag{16.1.3.1}
  \dim(S^{-1}A)\leq\dim(A).
\end{equation}
For every prime ideal $\mathfrak{p}$ of $A$, we have, by definition \sref{0.14.2.1}
\begin{equation}
\label{0.16.1.3.2}
\tag{16.1.3.2}
  \dim(A_\mathfrak{p})=\codim(V(\mathfrak{p}),\Spec(A)).
\end{equation}
This number is also called the \emph{height} of the prime ideal $\mathfrak{p}$ and denoted $\operatorname{ht}(\mathfrak{p})$.
More generally, for every ideal $\mathfrak{a}$ of $A$, we call the \emph{height} of $\mathfrak{a}$ the number
\begin{equation}
\label{0.16.1.3.3}
\tag{16.1.3.3}
  \operatorname{ht}(\mathfrak{a})=\codim(V(\mathfrak{a}),\Spec(A)).
\end{equation}
If $(\mathfrak{p}_\lambda)$ is the family of minimal prime ideals among those containing $\mathfrak{a}$, we have then \sref{0.14.2.1}
\begin{equation}
\label{0.16.1.3.4}
\tag{16.1.3.4}
  \operatorname{ht}(\mathfrak{a})=\inf_\lambda(\operatorname{ht}(\mathfrak{p}_\lambda)).
\end{equation}
We evidently have
\begin{equation}
\label{0.16.1.3.5}
\tag{16.1.3.5}
  \dim(A)=\sup_\mathfrak{m}(\dim(A_\mathfrak{m}))
\end{equation}
where $\mathfrak{m}$ ranges over the set of maximal ideals of $A$.
\end{env}

\begin{env}[16.1.4]
\label{0.16.1.4}
For every prime ideal $\mathfrak{p}$ of $A$, we have \sref{0.14.2.2.2}
\begin{equation}
\label{0.16.1.4.1}
\tag{16.1.4.1}
  \dim(A_\mathfrak{p})+\dim(A/\mathfrak{p})\leq\dim A.
\end{equation}
We say that $A$ is \emph{catenary} (resp.~\emph{equidimensional}, \emph{equicodimensional}, \emph{biequidimensional}) when $\Spec(A)$ is catenary (resp.~equidimensional, equicodimensional, biequidimensional).
When $A$ is Noetherian and biequidimensional, the two members of \sref{0.16.1.4.1} are equal for every prime ideal $\mathfrak{p}$ of $A$ \sref{0.14.3.5}.

For every ideal $\mathfrak{a}$ of $A$, the prime ideals of $A/\mathfrak{a}$ correspond bijectively to the prime ideals of $A$ containing $\mathfrak{a}$, hence, \emph{if $A$ is catenary, so is $A/\mathfrak{a}$}.
Similarly, for every multiplicative subset $S$ of $A$, the prime ideals of $S^{-1}A$ correspond bijectively to the prime ideals of $A$ not meeting $S$, hence \emph{if $A$ is catenary, so is $S^{-1}A$}.

Furthermore, as the prime ideals of $A$ contained in a prime ideal $\mathfrak{p}$ correspond bijectively to the prime ideals of $A_\mathfrak{p}$, for $A$ to be catenary it is necessary and sufficient that $A_\mathfrak{p}$ be so for every $\mathfrak{p}$.
To say that $A$ is catenary means (\sref{0.14.3.2} and \sref{0.16.1.3.2}) that for every pair of prime ideals $\mathfrak{p}$, $\mathfrak{q}$ of $A$ such that $\mathfrak{q}\subset\mathfrak{p}$, we have
\begin{equation}
\label{0.16.1.4.2}
\tag{16.1.4.2}
  \dim(A_\mathfrak{q})+\dim(A_\mathfrak{p}/\mathfrak{q}A_\mathfrak{p})=\dim(A_\mathfrak{p}).
\end{equation}
If $A$ is biequidimensional, it is the same for $A_\mathfrak{p}$ for every prime ideal $\mathfrak{p}$ of $A$.
\end{env}

\begin{proposition}[16.1.5]
\label{0.16.1.5}
Let $\rho:A\to B$ be a ring homomorphism making $B$ an integral $A$-algebra.
Then $\dim(B)\leq\dim(A)$; if in addition $\rho$ is injective, we have $\dim(B)=\dim(A)$.
\end{proposition}

\begin{proof}
\oldpage[0\textsubscript{IV-1}]{24}
If $\mathfrak{n}$ is the kernel of $\rho$, $\rho(A)$ is isomorphic to $A/\mathfrak{n}$, hence $\dim\rho(A)\leq\dim(A)$ by \sref{0.16.1.2.1}, and since $B$ is integral over $\rho(A)$, one can reduce to proving the second assertion.
If $\mathfrak{p}_0'\subset\mathfrak{p}_1'\subset\ldots\subset\mathfrak{p}_n'$ is a chain of prime ideals of $B$, the $\mathfrak{p}_i'\cap A$ are then distinct prime ideals in $A$ (Bourbaki, \emph{Alg.~comm.}, chap.~V, \textsection~2, n\textsuperscript{o}~1, cor.~1 of prop.~1), hence they form a chain of prime ideals.
Conversely, for every chain $\mathfrak{p}_0\subset\mathfrak{p}_1\subset\ldots\subset\mathfrak{p}_m$ of prime ideals of $A$, we know that for each $i$ there exists a prime ideal $\mathfrak{p}_i'$ of $B$ such that $\mathfrak{p}_i'\cap A=\mathfrak{p}_i$ and $\mathfrak{p}_i'\subset\mathfrak{p}_{i+1}'$ for $0\leq i\leq m\setmin 1$ (\emph{loc.~cit.}, cor.~2 of th.~1).
Hence the conclusion.
\end{proof}

\begin{proposition}[16.1.6]
\label{0.16.1.6}
Let $B$ be an integral ring, $A$ a local subring of $B$, $\mathfrak{m}$ its maximal ideal.
Suppose that $A$ is integrally closed and that $B$ is integral over $A$; then, for every maximal ideal $\mathfrak{n}$ of $B$, we have $\dim(B_\mathfrak{n})=\dim(A)$.
\end{proposition}

\begin{proof}
The first part of the reasoning of \sref{0.16.1.5} shows that $\dim(B_\mathfrak{n})\leq\dim(A)$.
Conversely, consider a chain of prime ideals $\mathfrak{p}_0\subset\mathfrak{p}_1\subset\ldots\subset\mathfrak{p}_k=\mathfrak{m}$ of $A$; we know that $\mathfrak{n}\cap A=\mathfrak{m}$ (Bourbaki, \emph{Alg.~comm.}, chap.~V, \textsection~2, n\textsuperscript{o}~1, prop.~1) and by virtue of the hypotheses made, one can apply the second theorem of Cohen-Seidenberg (Bourbaki, \emph{Alg.~comm.}, chap.~V, \textsection~2, n\textsuperscript{o}~4, th.~3), which proves the existence of a chain of prime ideals $\mathfrak{p}_0'\subset\mathfrak{p}_1'\subset\ldots\subset\mathfrak{p}_k'=\mathfrak{n}$ of $B$ such that $\mathfrak{p}_i'\cap A=\mathfrak{p}_i$ for each $i$; from this one deduces immediately $\dim(B_\mathfrak{n})\geq\dim(A)$, whence the proposition.
\end{proof}

\begin{env}[16.1.7]
\label{0.16.1.7}
Let $M$ be an $A$-module; we call the \emph{dimension of} $M$ and denote by $\dim_A(M)$ or $\dim(M)$ the dimension of the ring $A/\operatorname{Ann}(M)$ (isomorphic to the ring of homotheties $A_M$ of $M$).
We thus have $\dim(M^k)=\dim(M)$ for every $k>0$.
If $M$ is \emph{of finite type}, the prime ideals of $A$ containing $\operatorname{Ann}(M)$ are exactly those belonging to $\Supp(M)$ $(\mathbf{0}_\mathrm{I},~1.7.4)$ and the latter is closed in $\Spec(A)$, hence we then have
\begin{equation}
\label{0.16.1.7.1}
\tag{16.1.7.1}
  \dim(M)=\dim(\Supp(M))\leq\dim(A).
\end{equation}
If $M$, $N$ are two $A$-modules such that $N\subset M$, we have
\begin{equation}
\label{0.16.1.7.2}
\tag{16.1.7.2}
  \dim(N)\leq\dim(M),\qquad\dim(M/N)\leq\dim(M).
\end{equation}
If $S$ is a multiplicative subset of $A$ and if $M$ is of finite type, we have
\begin{equation}
\label{0.16.1.7.3}
\tag{16.1.7.3}
  \dim_{S^{-1}A}(S^{-1}M)\leq\dim_A(M).
\end{equation}
Furthermore, every maximal chain of prime ideals containing $\operatorname{Ann}(M)$ terminates at a maximal ideal $\mathfrak{m}$, hence its length is equal to the length of the chain of corresponding prime ideals of $A_\mathfrak{m}$ (which contain $(\operatorname{Ann}(M))_\mathfrak{m}$); from this remark and from \sref{0.16.1.7.3}, one obtains
\begin{equation}
\label{0.16.1.7.4}
\tag{16.1.7.4}
  \dim_A(M)=\sup_\mathfrak{m}(\dim_{A_\mathfrak{m}}(M_\mathfrak{m}))
\end{equation}
where $\mathfrak{m}$ ranges over the set of maximal ideals of $A$.
We say that $M$ is \emph{catenary} (resp.~\emph{equidimensional}, \emph{equicodimensional}, \emph{biequidimensional}) when $A/\operatorname{Ann}(M)$ is.
\end{env}

\begin{proposition}[16.1.8]
\label{0.16.1.8}
Let $A$ be a ring, $M$ an $A$-module of finite type, $\mathfrak{p}$ a prime ideal of $A$.
We have
\begin{equation}
\label{0.16.1.8.1}
\tag{16.1.8.1}
  \dim_{A_\mathfrak{p}}(M_\mathfrak{p})+\dim_A(M/\mathfrak{p}M)\leq\dim_A(M).
\end{equation}
\end{proposition}

\begin{proof}
\oldpage[0\textsubscript{IV-1}]{25}
We have $\operatorname{Ann}(M_\mathfrak{p})=(\operatorname{Ann}(M))_\mathfrak{p}$, hence the prime ideals of $A_\mathfrak{p}$ containing $\operatorname{Ann}(M_\mathfrak{p})$ correspond bijectively to the prime ideals of $A$ containing $\operatorname{Ann}(M)$ and contained in $\mathfrak{p}$.
On the other hand, if $\operatorname{Ann}(M)=\mathfrak{n}$, $V(\operatorname{Ann}(M/\mathfrak{p}M))=V(\mathfrak{p}+\mathfrak{n})$ $(\mathbf{0}_\mathrm{I},~1.7.5)$, hence the prime ideals of $A$ containing $\operatorname{Ann}(M/\mathfrak{p}M)$ contain $\mathfrak{p}+\mathfrak{n}$; whence the conclusion.
We note in addition that
\begin{equation}
\label{0.16.1.8.2}
\tag{16.1.8.2}
  \dim_{A/\mathfrak{p}}(M/\mathfrak{p}M)=\dim_A(M/\mathfrak{p}M).
\end{equation}
\end{proof}

\begin{proposition}[16.1.9]
\label{0.16.1.9}
Let $A$ be a Noetherian ring, $\rho:A\to B$ a ring homomorphism, $M$ a $B$-module such that $M_{[\rho]}$ is an $A$-module of finite type.
Then
\begin{equation}
\label{0.16.1.9.1}
\tag{16.1.9.1}
  \dim_A(M_{[\rho]})=\dim_B(M).
\end{equation}
\end{proposition}

\begin{proof}
Indeed, the ring $B_M$ of homotheties of the $B$-module $M$ is a subring of the $A$-module $C=\operatorname{End}_A(M_{[\rho]})$, and $C$ is an $A$-module of finite type (Bourbaki, \emph{Alg.~comm.}, chap.~III, \textsection~3, n\textsuperscript{o}~1, lemma~2), hence it is the same for $B_M$; moreover the ring $A_M$ of homotheties of the $A$-module $M_{[\rho]}$ is a subring of $B_M$, canonical image of $A$ in $C$, hence $B_M$ is finite over $A_M$; the conclusion then follows from \sref{0.16.1.5}.
\end{proof}

\begin{env}[16.1.10]
\label{0.16.1.10}
Let $A$ be a Noetherian ring; if $M$ is an $A$-module of \emph{finite length}, we know (Bourbaki, \emph{Alg.~comm.}, chap.~IV, \textsection~2, n\textsuperscript{o}~5, prop.~7 and cor.~1 of prop.~7) that $\Supp(M)$ is a finite discrete set, hence $\dim(M)=0$ if $M\neq 0$; conversely, if $M$ is an $A$-module of finite type such that $\dim(M)=0$, every point of $\Supp(M)$ is closed, in other words it is a maximal ideal of $A$, hence (\emph{loc.~cit.}, prop.~7) $M$ is of finite length.
\end{env}

\subsection[Dimension of a semi-local Noetherian ring]{Dimension of a semi-local Noetherian ring\protect\footnote{This section reproduces the lecture given by Serre in a course at the Collège de France in 1958.}.}
\label{subsection:0.16.2}

\begin{env}[16.2.1]
\label{0.16.2.1}
Let $A$ be a \emph{semi-local Noetherian} ring, $\mathfrak{r}$ its radical; recall that an \emph{ideal of definition} $\mathfrak{q}$ of $A$ is an ideal such that $\mathfrak{q}\subset\mathfrak{r}$ and $\mathfrak{q}$ contains a power of $\mathfrak{r}$, or equivalently, that the only prime ideals containing $\mathfrak{q}$ are \emph{maximal}, or that $A/\mathfrak{q}$ is an \emph{Artinian} ring.
Let $M$ be an $A$-module of finite type; for every integer $n>0$, $M/\mathfrak{q}^n M$ is an $A$-module of \emph{finite length}, equal to $\sum_{i=0}^{n\setmin 1}\operatorname{long}(\gr_i(M))$, where $\gr_\bullet(M)$ is the $(A/\mathfrak{q})$-graded module associated to $M$ equipped with the $\mathfrak{q}$-preadic filtration.
Recall that there exists a \emph{polynomial} $Q(n)$ with rational coefficients such that $\operatorname{long}(\gr_n(M))=Q(n)$ for $n$ sufficiently large $(\mathbf{III},~2.5.4)$; from this one deduces immediately that there exists a polynomial $P_\mathfrak{q}(M,n)$ in $n$, with rational coefficients, such that $\operatorname{long}(M/\mathfrak{q}^n M)=P_\mathfrak{q}(M,n)$ \emph{for $n$ sufficiently large}; it is clear that this polynomial is unique and that its leading coefficient is $>0$ if $M\neq 0$; we call it the \emph{Hilbert--Samuel polynomial} of $M$ for the $\mathfrak{q}$-preadic filtration.
\end{env}

\oldpage[0\textsubscript{IV-1}]{26}

\begin{lemma}[16.2.2.1]
\label{0.16.2.2.1}
Let $(M_n)$ be a $\mathfrak{q}$-good filtration of $M$ \textup{(Bourbaki, \emph{Alg.~comm.}, chap.~III, \textsection~3, n\textsuperscript{o}~1)}; then $\operatorname{long}(M/M_n)$ is equal, for $n$ sufficiently large, to a polynomial in $n$ whose degree and leading coefficient are the same as those of $P_\mathfrak{q}(M,n)$.
\end{lemma}

\begin{proof}
Indeed, there exists $n_0$ such that $M_{n+1}=\mathfrak{q}M_n$ for $n\geq n_0$, whence the inclusions $\mathfrak{q}^{n+n_0}M\subset M_{n+n_0}=\mathfrak{q}^n M_{n_0}\subset\mathfrak{q}^n M\subset M_n$ for $n$ sufficiently large, which implies
\[
  P_\mathfrak{q}(M,n+n_0)\geq\operatorname{long}(M/M_{n+n_0})\geq P_\mathfrak{q}(M,n)\geq\operatorname{long}(M/M_n)
\]
hence the conclusion.
\end{proof}

If $\mathfrak{q}'$ is a second ideal of definition, one has $\mathfrak{q}'\supset\mathfrak{q}^m$ for some integer $m$, hence $P_{\mathfrak{q}'}(M,n)\leq P_\mathfrak{q}(M,mn)$; the \emph{degree} of $P_\mathfrak{q}(M,n)$ therefore \emph{does not depend on the ideal of definition} $\mathfrak{q}$; we denote it $d(M)$.

\begin{lemma}[16.2.2.2]
\label{0.16.2.2.2}
Let $0\to M'\to M\to M''\to 0$ be a short exact sequence of $A$-modules of finite type.
Then the polynomial $P_\mathfrak{q}(M,n)\setmin P_\mathfrak{q}(M',n)\setmin P_\mathfrak{q}(M'',n)$ has degree $\leq d(M')\setmin 1\leq d(M)\setmin 1$.
\end{lemma}

\begin{proof}
Indeed, the filtration of the $M_n'=M'\cap\mathfrak{q}^n M$ is $\mathfrak{q}$-good (Bourbaki, \emph{Alg.~comm.}, chap.~III, \textsection~3, n\textsuperscript{o}~1, prop.~1); as $\operatorname{long}(M/\mathfrak{q}^n M)=\operatorname{long}(M''/\mathfrak{q}^n M'')+\operatorname{long}(M'/M_n')$, the lemma follows immediately from \sref{0.16.2.2.1} applied to $M'$.
\end{proof}

\begin{theorem}[16.2.3]
\label{0.16.2.3}
\textup{(Krull-Chevalley-Samuel).}
Let $A$ be a semi-local Noetherian ring, $\mathfrak{r}$ its radical, $M\neq 0$ an $A$-module of finite type, $d(M)$ the degree of the Hilbert--Samuel polynomial of $M$ (for any ideal of definition of $A$); let $s(M)$ denote the lower bound of the number of elements $x_1,\ldots,x_n$ of $\mathfrak{r}$ such that $M/(x_1M+\ldots+x_n M)$ is of finite length.
Then
\[
  d(M)=s(M)=\dim(M).
\]
\end{theorem}

\begin{lemma}[16.2.3.1]
\label{0.16.2.3.1}
For every $x\in\mathfrak{r}$, let ${}_x M$ be the submodule of $M$ of elements annihilated by $x$.
Then:
\begin{enumerate}
  \item[\rm{(i)}] $s(M)\leq s(M/x M)+1$.
  \item[\rm{(ii)}] Let $\mathfrak{p}_i$ be the prime ideals belonging to $\Supp(M)$ and such that
  \[
    \dim(A/\mathfrak{p}_i)=\dim(M)\quad(1\leq i\leq m).
  \]
\end{enumerate}
If $x\notin\mathfrak{p}_i$ for every $i$, we have $\dim(M/x M)\leq\dim(M)\setmin 1$.
\begin{enumerate}
  \item[\rm{(iii)}] If $\mathfrak{q}$ is an ideal of definition of $A$, the polynomial
  \[
    P_\mathfrak{q}({}_x M,n)\setmin P_\mathfrak{q}(M/x M,n)
  \]
  has degree $\leq d(M)\setmin 1$.
\end{enumerate}
\end{lemma}

\begin{proof}
The assertion (i) is trivial, since if $(x_i)_{1\leq i\leq m}$ is such that for the module $N=M/x M$, $N/(x_1N+\ldots+x_m N)$ is of finite length, this amounts to saying that $M/(x M+x_1M+\ldots+x_m M)$ is isomorphic to this module.

For (ii), note that, if $\mathfrak{a}$ denotes the annihilator of $M$, the prime ideals containing the annihilator of $M/x M$ are those containing $\mathfrak{a}+Ax$ $(\mathbf{0}_\mathrm{I},~1.7.5)$; none of the $\mathfrak{p}_i$ can therefore contain $\mathfrak{a}+Ax$ and by definition of $\dim(M)$ \sref{0.16.1.7} every chain of prime ideals containing $\mathfrak{a}+Ax$ has therefore length $\leq\dim(M)\setmin 1$.

\oldpage[0\textsubscript{IV-1}]{27}

Finally, to prove (iii), note that one has two short exact sequences
\[
  0\to{}_x M\to M\xrightarrow{x}x M\to 0,\qquad 0\to x M\to M\to M/x M\to 0
\]
and the assertion follows immediately from the lemma \sref{0.16.2.2.2}.
\end{proof}

The proof of \sref{0.16.2.3} proceeds in three steps:

\begin{env}[16.2.3.2]
\label{0.16.2.3.2}
$\dim(M)\leq d(M)$.
\end{env}

We argue by induction on $d(M)$, the case $d(M)=0$ being trivial.
Suppose thus $d(M)\geq 1$, and let $\mathfrak{p}_0\in\Supp(M)$ such that $\dim(A/\mathfrak{p}_0)=\dim(M)$.
This implies that $\mathfrak{p}_0$ is a \emph{minimal} element of $\Supp(M)$, hence is associated to $M$ (Bourbaki, \emph{Alg.~comm.}, chap.~IV, \textsection~1, n\textsuperscript{o}~4, th.~2) and $M$ therefore contains a submodule $N$ isomorphic to $A/\mathfrak{p}_0$ (\emph{loc.~cit.}, \textsection~1, n\textsuperscript{o}~1); as $d(M)\geq d(N)$ by \sref{0.16.2.2.2}, it suffices to prove that $\dim(N)\leq d(N)$.
Let then $\mathfrak{p}_0\subset\mathfrak{p}_1\subset\ldots\subset\mathfrak{p}_n$ be a chain of distinct prime ideals in $A$, and let us show that $n\leq d(N)$.
This is evident if $n=0$.
Otherwise, there exists $x\in\mathfrak{p}_1\cap\mathfrak{r}$ such that $x\notin\mathfrak{p}_0$, for the relation $\mathfrak{p}_0\supset\mathfrak{p}_1\cap\mathfrak{r}$ would imply $\mathfrak{p}_0\supset\mathfrak{r}$ since $\mathfrak{p}_0$ does not contain $\mathfrak{p}_1$, hence $\mathfrak{p}_0$ would be maximal, which is absurd.
We have $\mathfrak{p}_i\in\Supp(N/x N)$, for $i\geq 1$, hence $n\setmin 1\leq\dim(N/x N)$ \sref{0.16.1.7}.
As ${}_x N=0$ by the choice of $x$, we have $d(N/x N)\leq d(N)\setmin 1$ by \sref{0.16.2.3.1}, (iii).
The induction hypothesis then implies $d(N/x N)=\dim(N/x N)\geq n\setmin 1$, hence $n\setmin 1\leq d(N)\setmin 1$.

\begin{env}[16.2.3.3]
\label{0.16.2.3.3}
$d(M)\leq s(M)$.
\end{env}

Let indeed $(x_i)_{1\leq i\leq n}$ be a family of elements of $\mathfrak{r}$ such that, if we set $\mathfrak{a}=x_1A+\ldots+x_n A$, $M/\mathfrak{a}M$ is of finite length.
The ideals associated to $M/\mathfrak{a}M$ are then maximal (Bourbaki, \emph{Alg.~comm.}, chap.~IV, \textsection~2, n\textsuperscript{o}~5, prop.~7) and since these are the only prime ideals containing $\mathfrak{q}=\mathfrak{a}+(\mathfrak{r}\cap\operatorname{Ann}(M))$, $A/\mathfrak{q}$ is Artinian (\emph{loc.~cit.}, prop.~9) and $\mathfrak{q}$ is an ideal of definition of $A$.
But as $\mathfrak{q}^m M/\mathfrak{q}^{m+1}M=\mathfrak{a}^m M/\mathfrak{a}^{m+1}M$, and if $M$ is generated by $r$ elements $z_j$, $\mathfrak{a}^m M/\mathfrak{a}^{m+1}M$ is an $(A/\mathfrak{q})$-module generated by the canonical images of $x_1^{\nu_1}x_2^{\nu_2}\ldots x_n^{\nu_n}z_j$ where $\sum_i\nu_i=m$, hence by $r\binom{m+n\setmin 1}{n\setmin 1}$ elements; its length is therefore $\leq ar\binom{m+n\setmin 1}{n\setmin 1}$, where $a$ is a constant, and one deduces immediately that the degree of $P_\mathfrak{q}(M,n)$ is $\leq n$; whence $d(M)\leq s(M)$.

\begin{env}[16.2.3.4]
\label{0.16.2.3.4}
$s(M)\leq\dim(M)$.
\end{env}

Let us note first that $n=\dim(M)$ is finite by \sref{0.16.2.3.2}.
We argue by induction on $n$; the proposition is evident for $n=0$, since $M$ is then of finite length \sref{0.16.1.10}.
Suppose $n\geq 1$, and let $\mathfrak{p}_i$ $(1\leq i\leq m)$ be the prime ideals of $\Supp(M)$ such that $\dim(A/\mathfrak{p}_i)=n$; the definition of the $\mathfrak{p}_i$ shows that they are minimal elements of $\Supp(M)$, hence finite in number (Bourbaki, \emph{Alg.~comm.}, chap.~IV, \textsection~1, n\textsuperscript{o}~4, th.~2).
As $n\geq 1$, the $\mathfrak{p}_i$ are not maximal, and there therefore exists $x\in\mathfrak{r}$ such that $x\notin\mathfrak{p}_i$ for every $i$ (Bourbaki, \emph{Alg.~comm.}, chap.~II, \textsection~1, n\textsuperscript{o}~1, prop.~2).
By virtue of \sref{0.16.2.3.1}, we have $s(M)\leq s(M/x M)+1$ and $\dim(M)\geq\dim(M/x M)+1$; the induction hypothesis gives the conclusion.

\begin{corollary}[16.2.4]
\label{0.16.2.4}
Under the hypotheses of \sref{0.16.2.3}, we have
\begin{equation}
\label{0.16.2.4.1}
\tag{16.2.4.1}
  \dim_{\widehat{A}}(\widehat{M})=\dim_A(M).
\end{equation}
\end{corollary}

\oldpage[0\textsubscript{IV-1}]{28}

\begin{proof}
Indeed, $M/\widehat{\mathfrak{r}}^n M$ and $\widehat{M}/\mathfrak{r}^n\widehat{M}$ are isomorphic, hence $d(\widehat{M})=d(M)$.
\end{proof}

\begin{corollary}[16.2.5]
\label{0.16.2.5}
The dimension of a semi-local Noetherian ring $A$ is finite and equal to the minimum number of elements of the radical of $A$ generating an ideal of definition.
\end{corollary}

\begin{proof}
This is the equality $s(M)=\dim(M)$ for $M=A$.
\end{proof}

In particular:

\begin{corollary}[16.2.6]
\label{0.16.2.6}
Let $A$ be a local Noetherian ring, $\mathfrak{m}$ its maximal ideal, $k=A/\mathfrak{m}$ its residue field; we have
\begin{equation}
\label{0.16.2.6.1}
\tag{16.2.6.1}
  \dim(A)\leq\rg_k(\mathfrak{m}/\mathfrak{m}^2).
\end{equation}
\end{corollary}

\begin{proof}
Indeed, we know that if the $x_i$ $(1\leq i\leq n)$ are elements of $\mathfrak{m}$ whose classes mod.~$\mathfrak{m}^2$ form a basis of the $k$-vector space $\mathfrak{m}/\mathfrak{m}^2$, the $x_i$ generate $\mathfrak{m}$ (Bourbaki, \emph{Alg.~comm.}, chap.~II, \textsection~3, n\textsuperscript{o}~2, cor.~2 of prop.~4).
\end{proof}

\begin{proposition}[16.2.7]
\label{0.16.2.7}
Let $k$ be a field, $B$ a graded $k$-algebra of finite type generated by its elements of degree $1$ and such that $B_0=k$; let $E$ be a graded $B$-module of finite type, such that $(\mathbf{III},~2.5.4)$ for $n$ sufficiently large, $\rg_k(E_n)=r(n)$ is of the form $\Phi(n)\setmin\Phi(n\setmin 1)$, where $\Phi$ is a polynomial in $n$ of degree $d$.
Then there exist $d$ homogeneous elements of degree $\geq 1$ in $B$, $t_1,\ldots,t_d$, forming a family $(t_i)_{1\leq i\leq d}$ such that $E/(\sum_i t_i E)$ is of finite rank over $k$.
\end{proposition}

\begin{proof}
Let $\mathfrak{q}$ be the maximal ideal $\bigoplus_{n\geq 1}B_n$ of $B$, and, for every $n$, let $E_n'$ be the sub-$B$-module $\bigoplus_{h\geq n+1}E_h$ of $E$.
Every element of $B\setmin\mathfrak{q}$ determines an injective homothety in the Artinian $B$-module $E/E_n'$, and this homothety is thus bijective (Bourbaki, \emph{Alg.}, chap.~VIII, \textsection~2, n\textsuperscript{o}~2, lemma~3), which implies that $(E/E_n')_\mathfrak{q}$ identifies with $E/E_n'$; as the residue field of $B_\mathfrak{q}$ is $k$, one deduces from the hypothesis that $\operatorname{long}_{B_\mathfrak{q}}((E/E_n')_\mathfrak{q})$ is a polynomial in $n$ of degree $d$. Moreover, $E$ is a $B$-module of finite type and $B$ is generated over $k$ by its homogeneous elements of degree $1$, hence the filtration $(E_n')$ is $\mathfrak{q}$-good (Bourbaki, \emph{Alg.~comm.}, chap.~III, \textsection~1, n\textsuperscript{o}~3, lemma~1), hence the filtration $((E_n')_\mathfrak{q})$ on $E_\mathfrak{q}$ is $\mathfrak{q}B_\mathfrak{q}$-good.
Since $B_\mathfrak{q}$ is a local Noetherian ring, one concludes from \sref{0.16.2.3} and \sref{0.16.2.2.1} that $\dim(E_\mathfrak{q})=d$. We argue by induction on $d$, the assertion being evident if $d=0$.
Suppose $d\geq 1$, and observe that the minimal prime ideals $\mathfrak{p}_i$ $(1\leq i\leq r)$ in $\operatorname{Ass}(E)$ are \emph{graded} (Bourbaki, \emph{Alg.~comm.}, chap.~IV, \textsection~3, n\textsuperscript{o}~1, prop.~1), and that $\mathfrak{q}$ is distinct from the union of the $\mathfrak{p}_i$ since $d\geq 1$; there is therefore a \emph{homogeneous} element $t_1\in\mathfrak{q}$ not belonging to any of the $\mathfrak{p}_i$ (Bourbaki, \emph{Alg.~comm.}, chap.~III, \textsection~1, n\textsuperscript{o}~4, prop.~8).
As the prime ideals $(\mathfrak{p}_i)_\mathfrak{q}$ are the minimal elements of $\operatorname{Ass}(E_\mathfrak{q})$ (Bourbaki, \emph{Alg.~comm.}, chap.~IV, \textsection~1, n\textsuperscript{o}~2, prop.~5), we have $\dim(E_\mathfrak{q}/t_1E_\mathfrak{q})=d\setmin 1$ \sref{0.16.3.4}.
It suffices then to apply the induction hypothesis to the graded $B$-module $E/t_1E$ to obtain the conclusion.
\end{proof}

\subsection{Systems of parameters in a Noetherian local ring}
\label{subsection:0.16.3}

\begin{proposition}[16.3.1]
\label{0.16.3.1}
Let $A$ be a Noetherian ring, $\mathfrak{p}$ a prime ideal of $A$, and $n$ an integer.
The following conditions are equivalent:
\begin{enumerate}
  \item[\textup{a)}] $\operatorname{ht}(\mathfrak{p})=\dim(A_\mathfrak{p})\leq n$;
  \item[\textup{b)}] There exist $n$ elements $x_i$ of $A$ $(1\leq i\leq n)$ such that $\mathfrak{p}$ is minimal in the set of prime ideals containing the ideal $\mathfrak{a}$ generated by the $x_i$ \textup{(in other words,} $V(\mathfrak{p})$ is an irreducible component of $V(\mathfrak{a})$).
\end{enumerate}
\end{proposition}

\begin{proof}
Condition b) implies that $\mathfrak{a}A_\mathfrak{p}$ is an ideal of definition of $A_\mathfrak{p}$, whence a) by virtue of \sref{0.16.2.5}.
Conversely, if a) holds, there exists an ideal of definition $\mathfrak{b}$ of $A_\mathfrak{p}$ generated by $n$ elements $x_i/s$, with $s\in A\setmin\mathfrak{p}$.
By virtue of the one-to-one correspondence between prime ideals of $A_\mathfrak{p}$ and prime ideals of $A$ contained in $\mathfrak{p}$, the ideal $\mathfrak{a}$ of $A$ generated by the $x_i$ satisfies b).
\end{proof}

\oldpage[0\textsubscript{IV-1}]{29}
For $n=1$, one obtains the special case:

\begin{corollary}[16.3.2]
\label{0.16.3.2}
\textup{(«~Hauptidealsatz~»).}
For a prime ideal in a Noetherian ring to be of height $\leq 1$, it is necessary and sufficient that it be minimal in the set of prime ideals containing a suitable principal ideal.
\end{corollary}

\begin{corollary}[16.3.3]
\label{0.16.3.3}
\textup{(Artin-Tate).}
Let $A$ be an integral Noetherian ring.
The following conditions are equivalent:
\begin{enumerate}
  \item[\textup{a)}] $A$ is a semi-local ring of dimension $\leq 1$.
  \item[\textup{b)}] There exists $f\neq 0$ in $A$ such that $A_f$ is a field.
\end{enumerate}
\end{corollary}

\begin{proof}
The condition $a)$ is equivalent to saying that $A$ has only finitely many prime ideals $\mathfrak{m}_i\neq 0$, and that these ideals are all maximal (0 being a prime ideal).
As the product of the $\mathfrak{m}_i$ is not zero, since $A$ is integral, an element $f\neq 0$ of this product belongs to every maximal ideal, hence $(0)$ is the only prime ideal of the integral ring $A_f$, in other words $A_f$ is a field.
(One notes that this part of the argument does not use the fact that $A$ is Noetherian.)
Let us now prove that $b)$ implies $a)$; the hypothesis $b)$ means that every prime ideal $\mathfrak{p}\neq 0$ of $A$ contains $f$.
Let $\mathfrak{p}_i$ $(1\leq i\leq n)$ be the minimal prime ideals among those containing $f$ (which are finite in number since $A/f A$ is Noetherian); it suffices to prove that the $\mathfrak{p}_i$ are \emph{maximal} ideals, for then every nonzero prime ideal necessarily contains one of the $\mathfrak{p}_i$ and is therefore equal to it.
Suppose then that there exists a maximal ideal $\mathfrak{m}$ containing one of the $\mathfrak{p}_i$ and distinct from the $\mathfrak{p}_i$; $\mathfrak{m}$ is therefore not contained in the union of the $\mathfrak{p}_i$ (Bourbaki, \emph{Alg.~comm.}, chap.~II, \textsection~1, n\textsuperscript{o}~1, prop.~2), in other words there exists $g\in\mathfrak{m}$ not belonging to any of the $\mathfrak{p}_i$.
Let $\mathfrak{q}$ be a minimal prime ideal among those containing $g$; it follows from the «~Hauptidealsatz~» \sref{0.16.3.2} that $\mathfrak{q}$ is of height 1, hence $\neq 0$; it therefore contains by assumption $f$, hence also one of the $\mathfrak{p}_i$, and as $g$ does not belong to any of the $\mathfrak{p}_i$, $\mathfrak{q}$ is distinct from the $\mathfrak{p}_i$; it would thus be of height $\geq 2$, which is a contradiction.
\end{proof}

\begin{proposition}[16.3.4]
\label{0.16.3.4}
Let $A$ be a semi-local Noetherian ring, $\mathfrak{r}$ its radical, $M$ an $A$-module of finite type, $\mathfrak{p}_i$ the elements of $\Supp(M)$ such that $\dim(A/\mathfrak{p}_i)=\dim(M)$.
Then the $\mathfrak{p}_i$ are minimal elements of $\operatorname{Ass}(M)$, and $\dim(M/\mathfrak{p}_i M)=\dim(M)$.
For every $x\in\mathfrak{r}$, we have
\begin{equation}
\label{0.16.3.4.1}
\tag{16.3.4.1}
  \dim(M/x M)\geq\dim(M)\setmin 1.
\end{equation}
Furthermore, for the two members of \sref{0.16.3.4.1} to be equal, it is necessary and sufficient that $x$ not belong to any of the $\mathfrak{p}_i$.
\end{proposition}

\begin{proof}
The first assertion is immediate, since the $\mathfrak{p}_i$ are by definition the minimal elements of $\Supp(M)$ \sref{0.16.1.6}, and they are also the minimal elements of $\operatorname{Ass}(M)$ (Bourbaki, \emph{Alg.~comm.}, chap.~IV, \textsection~1, n\textsuperscript{o}~4, th.~2).
Furthermore, $\Supp(M/\mathfrak{p}_i M)$ is the set of prime ideals of $A$ containing $\mathfrak{p}_i$ $(\mathbf{0}_\mathrm{I},~1.7.5)$, hence $\dim(M/\mathfrak{p}_i M)=\dim(A/\mathfrak{p}_i)=\dim(M)$.

Set $N=M/x M$.
If $x_1,\ldots,x_n$ are elements of $\mathfrak{r}$ such that $M/(x M+x_1M+\ldots+x_n M)$ is of finite length, we have $s(M)\leq n+1$, hence $\dim(M)\leq n+1$ by \sref{0.16.2.3}, and as one can choose them so that $n=s(N)=\dim(N)$, we get \sref{0.16.3.4.1}.

The fact that the two members of \sref{0.16.3.4.1} are equal when $x$ does not belong to any of the $\mathfrak{p}_i$ follows from \sref{0.16.3.4.1} and \sref{0.16.2.3.1}.
Conversely, if \oldpage[0\textsubscript{IV-1}]{30} $\dim(M/x M)=\dim(M)\setmin 1$, none of the $\mathfrak{p}_i$ can belong to $\Supp(M/x M)$, and the prime ideals in the support of the latter are those containing $\operatorname{Ann}(M)+Ax$; since the $\mathfrak{p}_i$ contain $\operatorname{Ann}(M)$, they cannot contain $x$.
\end{proof}

\begin{corollary}[16.3.5]
\label{0.16.3.5}
With the notations of \sref{0.16.3.4}, for every ideal $\mathfrak{a}\subset\mathfrak{p}_i$, we have $\dim(M/\mathfrak{a}M)=\dim(M)$ and if one sets $N=M/\mathfrak{a}M$, the relation $\dim(M/x M)=\dim(M)\setmin 1$ implies $\dim(N/x N)=\dim(N)\setmin 1$.
\end{corollary}

\begin{proof}
Indeed, $\Supp(N)$ is the set of prime ideals containing $\mathfrak{a}+\operatorname{Ann}(M)$, hence $\mathfrak{p}_i$ belongs to $\Supp(N)$ and $\dim(N)\leq\dim(M)=\dim(A/\mathfrak{p}_i)$, hence $\dim(A/\mathfrak{p}_i)=\dim(N)$; furthermore the prime ideals $\mathfrak{q}\in\Supp(N)$ such that $\dim(A/\mathfrak{q})=\dim(N)$ are certain of the $\mathfrak{p}_j$; if $x$ does not belong to any of the $\mathfrak{p}_j$, we have then $\dim(N/x N)=\dim(N)\setmin 1$ by virtue of \sref{0.16.3.4}.
\end{proof}

\begin{definition}[16.3.6]
\label{0.16.3.6}
Let $A$ be a semi-local Noetherian ring, $\mathfrak{r}$ its radical, $M$ an $A$-module of finite type, and set $n=\dim(M)$.
A system of parameters for $M$ is any family $(x_i)_{1\leq i\leq n}$ of $n$ elements of $\mathfrak{r}$ such that $M/(x_1M+\ldots+x_n M)$ is of \emph{finite length}.
\end{definition}

We have seen \sref{0.16.2.3} that such systems always exist.
We have seen in the course of the proof of \sref{0.16.2.3.3} that if $\mathfrak{J}=\sum_{i=1}^n x_i A$ is such that $M/\mathfrak{J}M$ is of finite length, $\mathfrak{q}=\mathfrak{J}+\operatorname{Ann}(M)$ is an \emph{ideal of definition} of $A$, and conversely, if this is so, $M/\mathfrak{J}M$ is a module of finite type over the Artinian ring $A/\mathfrak{q}$, hence is of finite length.
It amounts to the same to say that $(x_i)_{1\leq i\leq n}$ is a system of parameters \emph{for} $M$ or \emph{for} $A/\operatorname{Ann}(M)$ (or equivalently that their images in $A/\operatorname{Ann}(M)$ form a system of parameters for this ring).
If $(x_i)_{1\leq i\leq n}$ is a system of parameters for $M$, so is $(x_i^k)$ for every integer $k\geq 1$, since one can reduce, by the preceding, to the case $M=A$, and if $\mathfrak{J}$ is an ideal of definition of $A$, $\sum_i x_i^k A$ contains $\mathfrak{J}^{kn}$, hence is also an ideal of definition.

\begin{proposition}[16.3.7]
\label{0.16.3.7}
Let $A$ be a semi-local Noetherian ring, $\mathfrak{r}$ its radical, $x_1,\ldots,x_k$ elements of $\mathfrak{r}$, $M$ an $A$-module of finite type.
Then
\begin{equation}
\label{0.16.3.7.1}
\tag{16.3.7.1}
  \dim(M/(x_1M+\ldots+x_k M))\geq\dim(M)\setmin k.
\end{equation}
Furthermore, for the two members of \sref{0.16.3.7.1} to be equal, it is necessary and sufficient that there exist elements $x_i$ $(k+1\leq i\leq n=\dim(M))$ of $\mathfrak{r}$ such that $(x_i)_{1\leq i\leq n}$ is a system of parameters for $M$.
\end{proposition}

\begin{proof}
The inequality \sref{0.16.3.7.1} follows from \sref{0.16.3.4.1} by induction on $k$.
If the two members of \sref{0.16.3.7.1} are equal, and if $x_{k+1},\ldots,x_n$ is a system of parameters for $M/(x_1M+\ldots+x_k M)$, it is clear that $M/(x_1M+\ldots+x_k M+x_{k+1}M+\ldots+x_n M)$ is of finite length, hence $(x_i)_{1\leq i\leq n}$ is a system of parameters for $M$; conversely, if there exist $x_i$ $(k+1\leq i\leq n)$ having this property and if one sets $N=M/(x_1M+\ldots+x_k M)$, it is clear that $N/(x_{k+1}N+\ldots+x_n N)$ is of finite length, hence $\dim(N)\leq n\setmin k$, and the two members are equal by virtue of \sref{0.16.3.7.1}.
\end{proof}

\oldpage[0\textsubscript{IV-1}]{31}

\begin{proposition}[16.3.8]
\label{0.16.3.8}
Let $A$ be a semi-local Noetherian ring, $X=\Spec(A)$ its spectrum, $\mathfrak{a}$ an ideal of $A$ distinct from $A$ such that $\codim(V(\mathfrak{a}),X)=r>0$.
Then there exist $r$ elements $x_1,\ldots,x_r$ of $\mathfrak{a}$, forming part of a system of parameters of $A$, such that
\[
  \codim(V(\mathfrak{a}),V(Ax_1+\ldots+Ax_r))=0.
\]
\end{proposition}

\begin{proof}
Let $\mathfrak{p}_i$ $(1\leq i\leq m)$ be the minimal prime ideals of the ring $A$, $\mathfrak{q}_j$ $(1\leq j\leq n)$ the minimal prime ideals among those containing $\mathfrak{a}$; we have \sref{0.14.2.1} $\codim(V(\mathfrak{a}),X)=\inf_j(\codim(V(\mathfrak{q}_j),X))=\inf_j(\dim(A_{\mathfrak{q}_j}))$ \sref{0.16.1.3.2}.
The hypothesis $r>0$ implies that $\mathfrak{a}$ is not contained in any of the $\mathfrak{p}_i$; hence there exists $x\in\mathfrak{a}$ not belonging to any of the $\mathfrak{p}_i$ (Bourbaki, \emph{Alg.~comm.}, chap.~II, \textsection~1, n\textsuperscript{o}~1, prop.~2).
One has as above $\codim(V(\mathfrak{a}),V(Ax))=\inf_j(\dim((A/Ax)_{\mathfrak{q}_j/Ax}))=\inf_j(\dim(A_{\mathfrak{q}_j}/A_{\mathfrak{q}_j}x))$; but as $x/1$ is not in any of the minimal prime ideals of $A_{\mathfrak{q}_j}$, we have $\dim(A_{\mathfrak{q}_j}/A_{\mathfrak{q}_j}x)=\dim(A_{\mathfrak{q}_j})-1$ \sref{0.16.3.4}, whence $\codim(V(\mathfrak{a}),V(Ax))=r-1$.
For the same reason, $\dim(A/Ax)=\dim(A)-1$.
We now argue by induction on $r$; if $r>1$, there exist $r-1$ elements $x_i'$ $(2\leq i\leq r)$ of $A'=A/Ax$, forming part of a system of parameters of this ring, belonging to $\mathfrak{a}/Ax$, and such that $\codim(V(\mathfrak{a}/Ax),V(A'x_2'+\ldots+A'x_r'))=0$.
Let $(x_i')_{r+1\leq i\leq s}$ be a family of elements of $A'$ such that $(x_i')_{2\leq i\leq s}$ is a system of parameters of $A'$; for every $i$ such that $2\leq i\leq s$, let $x_i$ be a representative of the class $x_i'$ in $A$, with $x_i\in\mathfrak{a}$ for $2\leq i\leq r$. It is immediate that $A/(Ax+Ax_2+\ldots+Ax_s)$ is of finite length, and as $\dim(A)=s$, we see that $x_1=x$ and the $x_i$ for $i\geq 2$ satisfy the conditions of the statement.
\end{proof}

\begin{proposition}[16.3.9]
\label{0.16.3.9}
Let $A$, $B$ be two local Noetherian rings, $\mathfrak{m}$ the maximal ideal of $A$, $k$ its residue field, $\varphi:A\to B$ a local homomorphism.
\begin{enumerate}
  \item[\rm{(i)}] We have
  \begin{equation}
  \label{0.16.3.9.1}
  \tag{16.3.9.1}
    \dim(B)\leq\dim(A)+\dim(B\otimes_Ak).
  \end{equation}
  \item[\rm{(ii)}] Let $M\neq 0$ be an $A$-module of finite type, $N\neq 0$ a $B$-module of finite type.
  We have
  \begin{equation}
  \label{0.16.3.9.2}
  \tag{16.3.9.2}
    \dim_B(M\otimes_AN)\leq\dim_A(M)+\dim_{B\otimes_Ak}(N\otimes_Ak).
  \end{equation}
\end{enumerate}
\end{proposition}

\begin{proof}
Let $m$ be the dimension of $A$, $(s_i)_{1\leq i\leq m}$ a system of parameters of $A$ \sref{0.16.3.6}, so that $\mathfrak{a}=\sum_i s_i A$, $A/\mathfrak{a}$ is an $A$-module of finite length.
Then the Artinian ring $A/\mathfrak{a}$ has nilpotent maximal ideal $\mathfrak{m}/\mathfrak{a}$; the ideal $\mathfrak{m}B/\mathfrak{a}B$ of $B/\mathfrak{a}B$, image of $(\mathfrak{m}/\mathfrak{a})\otimes_AB$, is therefore also nilpotent, and hence $\dim(B/\mathfrak{a}B)=\dim(B/\mathfrak{m}B)$ (the spectra of these two rings having the same underlying space).
On the other hand $B/\mathfrak{a}B=B/(s_1B+\ldots+s_m B)$, and the images of the $s_i$ in $B$ belong to the maximal ideal of $B$.
Hence by \sref{0.16.3.7}, we have
\[
  \dim(B)\leq m+\dim(B\otimes_Ak)
\]
which is \sref{0.16.3.9.1}.

To prove \sref{0.16.3.9.2}, note that if $\mathfrak{r}$ (resp.~$\mathfrak{s}$) is the annihilator of $M$ (resp.~$N$) \oldpage[0\textsubscript{IV-1}]{32} in $A$ (resp.~$B$), we have $\dim_A(M)=\dim(A/\mathfrak{r})$, $\dim_B(N)=\dim(B/\mathfrak{s})$ by definition \sref{0.16.1.7}; on the other hand, $\Supp(M\otimes_AN)$ is a closed subset of $\Spec(B)$ that identifies $(\mathbf{I},~9.1.13)$ to $\Supp(M)\times_{\Spec(A)}\Supp(N)=\Spec((B/\mathfrak{s})\otimes_A(A/\mathfrak{r}))=\Spec(B/(\mathfrak{s}+\mathfrak{r}B))$.
As $N\otimes_Ak$ is a $B$-module of finite type, we have $\dim_{B\otimes_Ak}(N\otimes_Ak)=\dim_B(N\otimes_Ak)$ \sref{0.16.1.9}, and $\Supp(N\otimes_Ak)$ is equal to $\Spec(B/(\mathfrak{s}+\mathfrak{m}B))$ by the same reasoning as above.
Applying \sref{0.16.3.9.1} to $A/\mathfrak{r}$ and to $B/(\mathfrak{s}+\mathfrak{r}B)$, we get
\[
  \dim(B/(\mathfrak{s}+\mathfrak{r}B))\leq\dim(A/\mathfrak{r})+\dim(B/(\mathfrak{s}+\mathfrak{m}B))
\]
which is nothing other than \sref{0.16.3.9.2}.
\end{proof}

\begin{corollary}[16.3.10]
\label{0.16.3.10}
Under the hypotheses of \sref{0.16.3.9}, if $\mathfrak{m}B$ is an ideal of definition of $B$, we have $\dim(B)\leq\dim(A)$; if in addition $A$ is integral and $\dim(A)=\dim(B)$, $\varphi$ is injective.
\end{corollary}

\begin{proof}
The first assertion follows from \sref{0.16.3.9} since then $\dim(B/\mathfrak{m}B)=0$.
One can moreover replace $A$ by $\varphi(A)$, hence $\dim(B)\leq\dim(\varphi(A))$; if $\mathfrak{a}=\Ker(\varphi)\neq 0$, and if $A$ is integral, we have $\dim(\varphi(A))=\dim(A/\mathfrak{a})<\dim(A)$ \sref{0.16.1.2.2}, hence one can have $\dim(A)=\dim(B)$ only if $\mathfrak{a}=0$.
\end{proof}

\subsection[Depth and codepth]{Depth and
codepth\protect\footnote{In the third part of Chapter~III, we develop the
notion of depth from the cohomological point of view and in a more general
framework.}}
\label{subsection:0.16.4}

\begin{proposition}[16.4.1]
\label{0.16.4.1}
Let $A$ be a local Noetherian ring, $\mathfrak{m}$ its maximal ideal, $M$ an $A$-module of finite type.
Then every $M$-regular sequence $(x_i)_{1\leq i\leq r}$ \sref{0.15.1.7} formed of elements of $\mathfrak{m}$, is part of a system of parameters for $M$.
\end{proposition}

\begin{proof}
We argue by induction on $r$; consider the $A$-module
\[
  N=M/(x_1M+\ldots+x_{r-1}M);
\]
by the induction hypothesis $x_1,\ldots,x_{r-1}$ form part of a system of parameters for $M$, hence $\dim(N)=\dim(M)-(r-1)$ \sref{0.16.3.7}.
As by hypothesis the homothety of ratio $x_r$ in $N$ is injective, $x_r$ does not belong to any of the prime ideals associated to $N$, hence \sref{0.16.3.4} $\dim(N/x_r N)=\dim(N)-1$, which can also be written
\[
  \dim(M/(x_1M+\ldots+x_r M))=\dim(M)-r;
\]
the conclusion then follows from \sref{0.16.3.7}.
\end{proof}

For an $M$-regular sequence formed of elements of $\mathfrak{m}$, $M\neq 0$, there are at most $\dim(M)$ elements.

\begin{lemma}[16.4.2]
\label{0.16.4.2}
Let $A$ be a local Noetherian ring, $\mathfrak{m}$ its maximal ideal, $k=A/\mathfrak{m}$ its residue field, $M$ an $A$-module.
For an $M$-regular sequence $(x_i)_{1\leq i\leq n}$ of elements of $\mathfrak{m}$ to be maximal, it is necessary and sufficient that $\Hom_A(k,M/(x_1M+\ldots+x_n M))\neq 0$.
\end{lemma}

\begin{proof}
To say that $(x_i)$ is maximal means that, for no $x\in\mathfrak{m}$, the homothety of ratio $x$ in $N=M/(x_1M+\ldots+x_n M)$ is injective, or equivalently that $\mathfrak{m}$ is contained in the union of the prime ideals associated to $N$ (Bourbaki, \emph{Alg.~comm.}, chap.~IV, \textsection~1, n\textsuperscript{o}~1, cor.~2 of prop.~2), hence equal to one of them, since it is a \oldpage[0\textsubscript{IV-1}]{33} maximal ideal (Bourbaki, \emph{Alg.~comm.}, chap.~II, \textsection~1, n\textsuperscript{o}~1, prop.~2); in other words, there is an element $z\in N$ whose annihilator is $\mathfrak{m}$, and hence the submodule $Az$ of $N$ is isomorphic to $A/\mathfrak{m}=k$; whence the lemma.
\end{proof}

\begin{lemma}[16.4.3]
\label{0.16.4.3}
Let $A$ be a local Noetherian ring, $\mathfrak{m}$ its maximal ideal, $k=A/\mathfrak{m}$ its residue field, $M$ an $A$-module, $(x_i)_{1\leq i\leq n}$ an $M$-regular sequence of elements of $\mathfrak{m}$.
Then the $A$-modules $\Hom_A(k,M/(x_1M+\ldots+x_n M))$ and $\Ext^n_A(k,M)$ are isomorphic, and for $n\geq 1$, they are also isomorphic to $\Ext^{n-1}_A(k,M/x_1M)$.
\end{lemma}

\begin{proof}
We argue by induction on $n$, the proposition being evident for $n=0$.
Set $N=M/x_1M$; the sequence $(x_i)_{2\leq i\leq n}$ is $N$-regular, hence
\[
  \Hom_A(k,M/(x_1M+\ldots+x_n M))=\Hom_A(k,N/(x_2N+\ldots+x_n N))
\]
is isomorphic to $\Ext^{n-1}_A(k,N)$ by the induction hypothesis.
Consider then the short exact sequence $0\to M\xrightarrow{x_1}M\to M/x_1M=N\to 0$; one deduces the exact sequence of Ext
\begin{equation}
\label{0.16.4.3.1}
\tag{16.4.3.1}
  \Ext^{n-1}_A(k,M)\to\Ext^{n-1}_A(k,N)\to\Ext^n_A(k,M)\xrightarrow{x_1}\Ext^n_A(k,M).
\end{equation}
By the induction hypothesis, $\Ext^{n-1}_A(k,M)$ is isomorphic to $\Hom_A(k,M/(x_1M+\ldots+x_{n-1}M))$, which is zero since $(x_i)_{1\leq i\leq n-1}$ is not a maximal $M$-regular sequence \sref{0.16.4.2}.
On the other hand, as $x_1\in\mathfrak{m}$, the homothety of ratio $x_1$ in the $A$-module $\Hom_A(k,T)$ is zero for every $A$-module $T$, hence it is the same for every $A$-module $\Ext^m_A(k,T)$; the assertions of the lemma then follow immediately from the exact sequence \sref{0.16.4.3.1}.
\end{proof}

\begin{corollary}[16.4.4]
\label{0.16.4.4}
Let $A$ be a local Noetherian ring, $\mathfrak{m}$ its maximal ideal, $k=A/\mathfrak{m}$ its residue field, $M$ an $A$-module of finite type $\neq 0$.
All maximal $M$-regular sequences of elements of $\mathfrak{m}$ have the same number of elements, which is the smallest integer $n$ such that $\Ext^n_A(k,M)\neq 0$.
\end{corollary}

\begin{definition}[16.4.5]
\label{0.16.4.5}
Let $A$ be a local Noetherian ring, $\mathfrak{m}$ its maximal ideal, $M$ an $A$-module of finite type.
We call the upper bound of the numbers of elements in $M$-regular sequences formed of elements of $\mathfrak{m}$ the depth of $M$, and denote it $\operatorname{prof}_AM$ or $\operatorname{prof}(M)$.
\end{definition}

We thus have $\operatorname{prof}(M)=+\infty$ if and only if $M=0$, and, by \sref{0.16.4.1},
\begin{equation}
\label{0.16.4.5.1}
\tag{16.4.5.1}
  \operatorname{prof}(M)\leq\dim(M)\quad\text{if}\quad M\neq 0.
\end{equation}

We evidently have $\operatorname{prof}(M^k)=\operatorname{prof}(M)$ for every $k>0$.

\begin{proposition}[16.4.6]
\label{0.16.4.6}
Let $A$ be a local Noetherian ring, $\mathfrak{m}$ its maximal ideal, $M$ an $A$-module of finite type.
\begin{enumerate}
  \item[\rm{(i)}]\label{0.16.4.6.i} For $\operatorname{prof}(M)=0$, it is necessary and sufficient that $\mathfrak{m}\in\operatorname{Ass}(M)$ \textup{(or equivalently that there exists a submodule of $M$ isomorphic to the residue field $k=A/\mathfrak{m}$ of $A$)}.
  \item[\rm{(ii)}] For every $M$-regular element $x$ belonging to the maximal ideal of $A$, we have
  \begin{equation}
  \label{0.16.4.6.1}
  \tag{16.4.6.1}
    \operatorname{prof}(M/x M)=\operatorname{prof}(M)-1.
  \end{equation}
  \item[\rm{(iii)}] If $M\neq 0$, we have
  \begin{equation}
  \label{0.16.4.6.2}
  \tag{16.4.6.2}
    \operatorname{prof}(M)\leq\inf_{\mathfrak{p}\in\operatorname{Ass}(M)}\dim(A/\mathfrak{p})\leq\dim(M).
  \end{equation}
  \oldpage[0\textsubscript{IV-1}]{34}
  \item[\rm{(iv)}] We have $\operatorname{prof}_A(M)=\operatorname{prof}_{\widehat{A}}(\widehat{M})$.
\end{enumerate}
\end{proposition}

\begin{proof}
(i) To say $\operatorname{prof}(M)=0$ means that there is no $x\in\mathfrak{m}$ such that the homothety of ratio $x$ in $M$ is injective, that is, every $x\in\mathfrak{m}$ belongs to a prime ideal associated to $M$ (Bourbaki, \emph{Alg.~comm.}, chap.~IV, \textsection~1, n\textsuperscript{o}~1, cor.~2 of prop.~2); but $\mathfrak{m}$ cannot be the union of prime ideals associated to $M$ unless it is equal to one of them (Bourbaki, \emph{Alg.~comm.}, chap.~II, \textsection~1, n\textsuperscript{o}~1, prop.~2).

(ii) If $(x_i)_{1\leq i\leq m}$ is a $(M/x M)$-regular sequence of elements of $\mathfrak{m}$, the sequence formed of $x$ and the $x_i$ is $M$-regular, and it follows from the definition that it is maximal when the first sequence is; whence the conclusion.

(iii) We argue by induction on $n=\operatorname{prof}(M)$, the assertion being evident for $n=0$.
We use the following lemma:
\end{proof}

\begin{lemma}[16.4.6.3]
\label{0.16.4.6.3}
Let $A$ be a Noetherian ring, $M$ an $A$-module of finite type, $t\in A$ an $M$-regular element.
Then, for every $\mathfrak{p}\in\operatorname{Ass}(M)$, every \emph{minimal} prime ideal among those containing $\mathfrak{p}+At$ belongs to $\operatorname{Ass}(M/t M)$.
\end{lemma}

\begin{proof}[Proof of the lemma]
We know that there exists a short exact sequence
\[
  0\to M'\to M\to M''\to 0
\]
such that $\operatorname{Ass}(M')=\{\mathfrak{p}\}$, $\operatorname{Ass}(M'')=\operatorname{Ass}(M)\setmin\{\mathfrak{p}\}$ (Bourbaki, \emph{Alg.~comm.}, chap.~IV, \textsection~1, n\textsuperscript{o}~1, prop.~4); as $t$ is $M$-regular, it is also $M'$-regular and $M''$-regular (\emph{loc.~cit.}, \textsection~1, n\textsuperscript{o}~1, cor.~2 of prop.~2).
One deduces that the sequence
\[
  0\to M'/t M'\to M/t M\to M''/t M''\to 0
\]
is exact \sref{0.15.1.18}.
Consequently, $\operatorname{Ass}(M'/t M')\subset\operatorname{Ass}(M/t M)$.
But the prime ideals of $A$ containing $\mathfrak{p}$ are the points of $\Supp(M')$ (\emph{loc.~cit.}, \textsection~1, n\textsuperscript{o}~3, prop.~7), hence the prime ideals containing $\mathfrak{p}+At$ are the points of the support of $M'/t M'=M'\otimes_A(A/t A)$ $(\mathbf{0}_\mathrm{I},~1.7.5)$; a minimal element of the set of these prime ideals belongs then to $\operatorname{Ass}(M'/t M')$ (Bourbaki, \emph{loc.~cit.}, \textsection~1, n\textsuperscript{o}~3, cor.~1 of prop.~7), and \emph{a fortiori} to $\operatorname{Ass}(M/t M)$.
\end{proof}

\begin{proof}[Continuation of the proof of \sref{0.16.4.6}]
This lemma established, we return to the proof of (iii).
Let $x\in\mathfrak{m}$ be an $M$-regular element, so that $\operatorname{prof}(M/x M)=n-1$ by (ii); for every $\mathfrak{p}\in\operatorname{Ass}(M)$, it follows from the lemma that there is a prime ideal $\mathfrak{p}'$ associated to $M/x M$ and containing $\mathfrak{p}+x A$; but as $x$ is $M$-regular, we have $x\notin\mathfrak{p}$, whence $\mathfrak{p}'\neq\mathfrak{p}$ and hence $\dim(A/\mathfrak{p}')\leq\dim(A/\mathfrak{p})-1$.
The induction hypothesis shows that $n-1\leq\dim(A/\mathfrak{p}')$; one concludes that $n\leq\dim(A/\mathfrak{p})$ for every $\mathfrak{p}\in\operatorname{Ass}(M)$.

(iv) One can restrict to the case $M\neq 0$.
Recall that $\widehat{M}=M\otimes_A\widehat{A}$ up to canonical isomorphism $(\mathbf{0}_\mathrm{I},~7.3.3)$; as $\widehat{A}$ is a flat $A$-module, there are canonical isomorphisms $\Ext^i_A(k,M)\otimes_A\widehat{A}\xrightarrow{\sim}\Ext^i_{\widehat{A}}(k,M\otimes_A\widehat{A})$ since $k=k\otimes_A\widehat{A}$ up to canonical isomorphism $(\mathbf{0}_\mathrm{III},~12.3.5)$.
The assertion (iv) then follows from \sref{0.16.4.4} and the fact that $\widehat{A}$ is a faithfully flat $A$-module $(\mathbf{0}_\mathrm{I},~7.3.5)$.
\end{proof}

\begin{env}[16.4.6.4]
\label{0.16.4.6.4}
One notes that if $\mathfrak{p}$ is a prime ideal of $A$, one can \emph{also} have $\operatorname{prof}_{A_\mathfrak{p}}(M_\mathfrak{p})<\operatorname{prof}_A(M)$ and $\operatorname{prof}_{A_\mathfrak{p}}(M_\mathfrak{p})>\operatorname{prof}_A(M)$; for an example of the first case,
\oldpage[0\textsubscript{IV-1}]{35}
it suffices to take $A$ a local integral ring of dimension $\geq 1$; we have $\operatorname{prof}(A)\geq 1$ but $\operatorname{prof}(K)=0$ for the field of fractions $K$ of $A$.
For an example of the second case, consider a Noetherian integral local ring $A_0$ of dimension $\geq 2$, with maximal ideal $\mathfrak{m}_0$ and residue field $k=A_0/\mathfrak{m}_0$; let $A=A_0\oplus\mathfrak{J}$ be the trivial extension of $A_0$ by the $A_0$-module $k$ (so that the multiplication in $A$ is given by $(x,z)(x',z')=(xx',xz'+x'z)$).
It is clear that $\mathfrak{J}$ is the nilradical of $A$ and $A/\mathfrak{J}$ is isomorphic to $A_0$, so that every prime ideal $\mathfrak{p}$ of $A$ is of the form $\mathfrak{p}_0\oplus\mathfrak{J}$, where $\mathfrak{p}_0$ is a prime ideal of $A_0$; in particular $\mathfrak{m}=\mathfrak{m}_0\oplus\mathfrak{J}$ is the only maximal ideal of $A$.
As every element of $\mathfrak{m}_0$ annihilates $\mathfrak{J}$, we have $\operatorname{prof}(A)=\operatorname{prof}(A_\mathfrak{m})=0$; on the contrary, if $\mathfrak{p}_0\neq\mathfrak{m}_0$, the elements of $\mathfrak{J}$ are annihilated by those of $\mathfrak{m}_0\setmin\mathfrak{p}_0$, hence $A_\mathfrak{p}$ is canonically isomorphic to $(A_0)_{\mathfrak{p}_0}$, and as $A_0$ is integral, $\operatorname{prof}((A_0)_{\mathfrak{p}_0})\geq 1$ if $\mathfrak{p}_0\neq 0$.
\end{env}

\begin{corollary}[16.4.7]
\label{0.16.4.7}
For a local Noetherian reduced ring $A$ to be of depth $0$, it is necessary and sufficient that $A$ be a field.
\end{corollary}

\begin{proof}
Indeed, as the intersection of the minimal prime ideals of $A$ is $0$, $A$ has no embedded associated prime ideals; by \sref{0.16.4.6}, (i), if $\operatorname{prof}(A)=0$, the maximal ideal $\mathfrak{m}$ of $A$ is also a minimal prime ideal, hence it is the only prime ideal of $A$, and as $A$ is reduced, $\mathfrak{m}=0$.
\end{proof}

\begin{proposition}[16.4.8]
\label{0.16.4.8}
Let $A$, $B$ be two local Noetherian rings, $\rho:A\to B$ a local homomorphism, $M$ a $B$-module such that $M_{[\rho]}$ is an $A$-module of finite type.
Then
\begin{equation}
\label{0.16.4.8.1}
\tag{16.4.8.1}
  \operatorname{prof}_A(M_{[\rho]})=\operatorname{prof}_B(M).
\end{equation}
\end{proposition}

\begin{proof}
One can restrict to the case $M\neq 0$; suppose $\operatorname{prof}_A(M_{[\rho]})=n$.
If $(x_i)_{1\leq i\leq n}$ is a $(M_{[\rho]})$-regular maximal sequence formed of elements of the maximal ideal of $A$, the $\rho(x_i)$ constitute an $M$-regular sequence formed of elements of the maximal ideal of $B$, and we have then, setting $N=M/(\rho(x_1)M+\ldots+\rho(x_n)M)$, $\operatorname{prof}_B(N)=\operatorname{prof}_B(M)-n$ \sref{0.16.4.6}, (ii); similarly $N_{[\rho]}=M_{[\rho]}/(x_1M_{[\rho]}+\ldots+x_n M_{[\rho]})$ and $\operatorname{prof}_A(N_{[\rho]})=0$; one is thus reduced to proving the proposition when $n=0$.
Let $k$ be the residue field of $A$; since $k$ is a quotient of $A$, the $A$-module $P=\operatorname{Hom}_A(k,M_{[\rho]})$ is a submodule of $\operatorname{Hom}_A(A,M_{[\rho]})=M_{[\rho]}$, and the hypothesis implies that $P\neq 0$ \sref{0.16.4.6}, (i); furthermore, $P$ is an $A$-module of finite type (since $M_{[\rho]}$ is an $A$-module of finite type), hence a finite-dimensional $k$-vector space, annihilated by the maximal ideal of $A$, and finally an $A$-module \emph{of finite length}.
Moreover, $P$ is also a submodule of the \emph{$B$-module} $M$, and since it is of finite length as an $A$-module, it is \emph{a fortiori} of finite length as a $B$-module.
As it is $\neq 0$, it contains a simple sub-$B$-module, that is to say one isomorphic to the residue field of $B$; one concludes then by \sref{0.16.4.6}, (i) that $\operatorname{prof}_B(M)=0$.
\end{proof}

\begin{definition}[16.4.9]
\label{0.16.4.9}
Let $A$ be a local Noetherian ring, $M$ an $A$-module of finite type.
If $M\neq 0$, we call the codepth of $M$ and denote $\operatorname{coprof}_A(M)$ or $\operatorname{coprof}(M)$, the finite integer $\dim_A(M)-\operatorname{prof}_A(M)\geq 0$ \sref{0.16.4.5.1}.
If $M=0$, we set $\operatorname{coprof}(M)=0$.
\end{definition}

We evidently have $\operatorname{coprof}(M^k)=\operatorname{coprof}(M)$ for every $k>0$.

\begin{proposition}[16.4.10]
\label{0.16.4.10}
Let $A$ be a local Noetherian ring, $M$ an $A$-module of finite type.
\oldpage[0\textsubscript{IV-1}]{36}
\begin{enumerate}
  \item[\rm{(i)}] For every $M$-regular element $x$ belonging to the maximal ideal of $A$, we have $\operatorname{coprof}(M/x M)=\operatorname{coprof}(M)$.
  \item[\rm{(ii)}] We have $\operatorname{coprof}_A(M)=\operatorname{coprof}_{\widehat{A}}(\widehat{M})$.
\end{enumerate}
\end{proposition}

\begin{proof}
Indeed, (i) follows from \sref{0.16.3.4} and \sref{0.16.4.6}, (ii), and (ii) from \sref{0.16.2.4} and \sref{0.16.4.6}, (iv).
\end{proof}

We shall show later $(\mathbf{IV},~6.11.5)$ that for every prime ideal $\mathfrak{p}$ of $A$, we have $\operatorname{coprof}_{A_\mathfrak{p}}(M_\mathfrak{p})\leq\operatorname{coprof}_A(M)$.

\begin{proposition}[16.4.11]
\label{0.16.4.11}
Under the hypotheses of \sref{0.16.4.8}, we have
\begin{equation}
\label{0.16.4.11.1}
\tag{16.4.11.1}
  \operatorname{coprof}_A(M_{[\rho]})=\operatorname{coprof}_B(M).
\end{equation}
\end{proposition}

\begin{proof}
This follows from \sref{0.16.1.9} and \sref{0.16.4.8}.
\end{proof}

\subsection{Cohen-Macaulay modules}
\label{subsection:0.16.5}

\begin{definition}[16.5.1]
\label{0.16.5.1}
Let $A$ be a local Noetherian ring.
An $A$-module of finite type $M$ is called a Cohen-Macaulay module if $\operatorname{coprof}(M)=0$, that is, if $M=0$ or $\dim(M)=\operatorname{prof}(M)$.
We say that $A$ is a Cohen-Macaulay ring if the $A$-module $A$ is a Cohen-Macaulay module.
\end{definition}

An $A$-module of finite length is a Cohen-Macaulay module since it is of dimension 0 \sref{0.16.1.10}.
A local \emph{reduced} ring \emph{of dimension} 1 is a Cohen-Macaulay ring, since then its maximal ideal $\mathfrak{m}$ cannot be associated to $A$ (Bourbaki, \emph{Alg.~comm.}, chap.~IV, \textsection~2, n\textsuperscript{o}~5, prop.~10), hence $\operatorname{prof}(A)=1$ by virtue of \sref{0.16.4.6}, (i).
If $A$ is an \emph{integral local integrally closed} ring of dimension $\geq 2$, we have $\operatorname{prof}(A)\geq 2$: indeed, if $x\neq 0$ is an element of $\mathfrak{m}$, we know (Bourbaki, \emph{Alg.~comm.}, chap.~VII, \textsection~1, n\textsuperscript{o}~4, prop.~8) that the prime ideals associated to $A/x A$ are of height 1, hence distinct from $\mathfrak{m}$ and their union is therefore distinct from $\mathfrak{m}$ (Bourbaki, \emph{Alg.~comm.}, chap.~II, \textsection~1, n\textsuperscript{o}~1, prop.~2); there thus exist $(x,y)$ sequences that are $A$-regular and formed of elements of $\mathfrak{m}$, whence our assertion.
One concludes that an integral local \emph{integrally closed} ring \emph{of dimension} 2 is a Cohen-Macaulay ring.

\begin{proposition}[16.5.2]
\label{0.16.5.2}
Let $A$ be a local Noetherian ring, $M$ an $A$-module of finite type.
For $M$ to be a Cohen-Macaulay module, it is necessary and sufficient that $\widehat{M}$ be a $\widehat{A}$-module of Cohen-Macaulay.
\end{proposition}

\begin{proof}
This follows immediately from \sref{0.16.4.10}, (ii).
\end{proof}

\begin{proposition}[16.5.3]
\label{0.16.5.3}
Under the hypotheses of \sref{0.16.4.8}, for $M$ to be a Cohen-Macaulay $B$-module, it is necessary and sufficient that $M_{[\rho]}$ be a Cohen-Macaulay $A$-module.
\end{proposition}

\begin{proof}
This follows from \sref{0.16.4.11}.
\end{proof}

\begin{proposition}[16.5.4]
\label{0.16.5.4}
Let $A$ be a local Noetherian ring, $M$ an $A$-module of finite type $\neq 0$.
If $M$ is a Cohen-Macaulay module, we have $\dim(A/\mathfrak{p})=\dim(M)$ for every prime ideal $\mathfrak{p}\in\operatorname{Ass}(M)$; in particular no prime ideal associated to $M$ is embedded.
\end{proposition}

\begin{proof}
This follows immediately from the inequalities \sref{0.16.4.6.2}.
\end{proof}

One can also say that $M$ (or $\Supp(M)$) is \emph{equidimensional} \sref{0.16.1.7}.

\oldpage[0\textsubscript{IV-1}]{37}

\begin{proposition}[16.5.5]
\label{0.16.5.5}
Let $A$ be a local Noetherian ring, $\mathfrak{m}$ its maximal ideal, $M$ an $A$-module of finite type $\neq 0$, $x$ an element of $\mathfrak{m}$.
Suppose that $M$ is a Cohen-Macaulay module; then the following conditions are equivalent:
\begin{enumerate}
  \item[\textup{a)}] $x$ is $M$-regular;
  \item[\textup{b)}] $\dim(M/x M)=\dim(M)\setmin 1$;
  \item[\textup{c)}] $x$ does not belong to any \emph{minimal} element of $\Supp(M)$.
\end{enumerate}
Furthermore, $M/x M$ is then a Cohen-Macaulay module.
\end{proposition}

\begin{proof}
We have seen \sref{0.16.5.4} that all the $\mathfrak{p}$ in $\operatorname{Ass}(M)$ are minimal in $\Supp(M)$ and such that $\dim(A/\mathfrak{p})=\dim(M)$; the equivalence of a) and c) then follows from Bourbaki, \emph{Alg.~comm.}, chap.~IV, \textsection~1, n\textsuperscript{o}~1, cor.~2 of prop.~2; the equivalence of b) and c) follows from \sref{0.16.3.4}.
The fact that $M/x M$ is a Cohen-Macaulay module follows finally from \sref{0.16.4.10}, (i).
\end{proof}

\begin{corollary}[16.5.6]
\label{0.16.5.6}
Under the hypotheses of \sref{0.16.5.5}, let $(x_i)_{1\leq i\leq r}$ be a sequence of elements of $\mathfrak{m}$.
The following conditions are equivalent:
\begin{enumerate}
  \item[\textup{a)}] The sequence $(x_i)$ is $M$-regular.
  \item[\textup{b)}] $\dim(M/(x_1M+\ldots+x_r M))=\dim(M)\setmin r$.
  \item[\textup{c)}] The $x_i$ form part of a system of parameters for $M$.
\end{enumerate}
Furthermore, when these conditions are satisfied, $N=M/(x_1M+\ldots+x_r M)$ is a Cohen-Macaulay module, hence for every prime ideal $\mathfrak{p}\in\operatorname{Ass}(N)$, we have $\dim(A/\mathfrak{p})=\dim(M)\setmin r$.
\end{corollary}

\begin{proof}
We already know that \emph{b)} and \emph{c)} are equivalent \sref{0.16.3.6} and that \emph{a)} implies \emph{c)} \sref{0.16.4.1} without any hypothesis on $M$. To see that \emph{c)} implies \emph{a)}, and to prove the last assertion of the statement, we argue by induction on $r$, the assertion being trivial for $r=0$.
Since $x_1$ forms part of a system of parameters, $\dim(M/x_1M)=\dim(M)\setmin 1$ \sref{0.16.3.6}, hence $x_1$ is $M$-regular and $M/x_1M$ is Cohen-Macaulay by \sref{0.16.5.5}.
Since the sequence $(x_i)_{2\leq i\leq r}$ forms part of a system of parameters for $P=M/x_1M$, the induction hypothesis shows that $(x_i)_{2\leq i\leq r}$ is a $P$-regular sequence and that $P/(x_2P+\ldots+x_r P)=N$ is a Cohen-Macaulay module; furthermore $(x_i)_{1\leq i\leq r}$ is $M$-regular.
\end{proof}

\begin{remark}[16.5.7]
\label{0.16.5.7}
It follows from \sref{0.16.5.6} that there is an identity between \emph{systems of parameters for $M$ and the maximal $M$-regular sequences of elements of $\mathfrak{m}$} when $M$ is a Cohen-Macaulay module.
Conversely, it follows from \sref{0.16.4.1} that if there is \emph{any} $M$-regular sequence that is a system of parameters for $M$, then $M$ is a Cohen-Macaulay module.
The last property stated in \sref{0.16.5.6} also \emph{characterizes} Cohen-Macaulay modules:
\end{remark}

\begin{proposition}[16.5.8]
\label{0.16.5.8}
Let $A$ be a local Noetherian ring, $M$ an $A$-module of finite type.
Suppose that for every sequence $(x_i)_{1\leq i\leq r}$ of elements of $A$ forming part of a system of parameters for $M$, and for every prime ideal $\mathfrak{p}$ associated to $N=M/(x_1M+\ldots+x_r M)$, we have $\dim(A/\mathfrak{p})=\dim(N)$.
Then $M$ is a Cohen-Macaulay module.
\end{proposition}

\begin{proof}
Let $(y_i)_{1\leq i\leq n}$ be a system of parameters for $M$; it suffices to show that the sequence $(y_i)$ is $M$-regular.
We argue by induction on $n=\dim(M)$, the proposition being trivial for $n=0$.
By hypothesis (applied with $r=0$) the prime ideals $\mathfrak{p}_j$ associated to $M$ are all such that $\dim(A/\mathfrak{p}_j)=\dim(M)$; as $y_1$ forms part of a system of parameters
\oldpage[0\textsubscript{IV-1}]{38}
for $M$, it follows from \sref{0.16.3.7} and \sref{0.16.3.4} that $y_1$ does not belong to any of the $\mathfrak{p}_j$, hence is $M$-regular.
If we set $P=M/y_1M$, then $(y_i)_{2\leq i\leq n}$ is a system of parameters for $P$, and it is immediate that $P$ satisfies the hypothesis of the statement; applying the induction hypothesis, one sees that $(y_i)_{2\leq i\leq n}$ is a $P$-regular sequence, which proves the proposition.
\end{proof}

\begin{proposition}[16.5.9]
\label{0.16.5.9}
Let $A$ be a local Noetherian ring, $M$ an $A$-module of finite type.
Suppose that $M$ is a Cohen-Macaulay module.
Let $\mathfrak{p}\in\Supp(M)$, and set $r=\dim(M)\setmin\dim(M/\mathfrak{p}M)$.
Then there exists an $M$-regular sequence $(x_i)_{1\leq i\leq r}$, formed of elements of $\mathfrak{p}$; for every such sequence, we have
\[
  \dim(M/(x_1M+\ldots+x_r M))=\dim(M/\mathfrak{p}M)=\dim(A/\mathfrak{p}),
\]
and $\mathfrak{p}$ is a \emph{minimal} element of $\operatorname{Ass}(M/(x_1M+\ldots+x_r M))$.
\end{proposition}

\begin{proof}
Let us first prove the existence of the $x_i$; there is nothing to prove for $r=0$.
For $r>0$, we proceed by induction.
As for every prime ideal $\mathfrak{q}_i$ associated to $M$ we have $\dim(A/\mathfrak{q}_i)=\dim(M)$ \sref{0.16.5.4}, the hypothesis $\dim(M/\mathfrak{p}M)<\dim(M)$ implies that $\mathfrak{p}$ cannot be contained in any of the $\mathfrak{q}_i$, the latter being the minimal prime ideals among those containing $\operatorname{Ann}(M)$ \sref{0.16.1.2.2} and \sref{0.16.1.7}; one concludes that $\mathfrak{p}$ is not contained in the union of the $\mathfrak{q}_i$ (Bourbaki, \emph{Alg.~comm.}, chap.~II, \textsection~1, n\textsuperscript{o}~1, prop.~2), and hence there exists $x_1\in\mathfrak{p}$ that is $M$-regular (Bourbaki, \emph{Alg.~comm.}, chap.~IV, \textsection~1, n\textsuperscript{o}~1, cor.~2 of prop.~2).
It follows from \sref{0.16.5.5} that $N=M/x_1M$ is a Cohen-Macaulay module, of dimension $\dim(M)\setmin 1$; as $x_1\in\mathfrak{p}$, moreover $N/\mathfrak{p}N=M/\mathfrak{p}M$, whence $\dim(N)\setmin\dim(N/\mathfrak{p}N)=r\setmin 1$; one can therefore apply the induction hypothesis to $N$, and if $(x_i)_{2\leq i\leq r}$ is an $N$-regular sequence formed of elements of $\mathfrak{p}$, it is clear that $(x_i)_{1\leq i\leq r}$ is an $M$-regular sequence formed of elements of $\mathfrak{p}$.

Now if one sets $P=M/(x_1M+\ldots+x_r M)$, we have $P/\mathfrak{p}P=M/\mathfrak{p}M$ since the $x_i$ are in $\mathfrak{p}$, and $\dim(P)=\dim(P/\mathfrak{p}P)$ by \sref{0.16.5.6}.
But since $P$ is a Cohen-Macaulay module, we also have $\dim(P)=\dim(A/\mathfrak{p}')$ for every $\mathfrak{p}'\in\operatorname{Ass}(P)$ \sref{0.16.5.4}.
As $\mathfrak{p}\in\Supp(P)$ $(\mathbf{0}_\mathrm{I},~1.7.5)$, $\mathfrak{p}$ contains some $\mathfrak{p}'\in\operatorname{Ass}(P)$ (Bourbaki, \emph{Alg.~comm.}, chap.~IV, \textsection~1, n\textsuperscript{o}~4, th.~2), and as $\dim(P)=\dim(P/\mathfrak{p}P)\leq\dim(A/\mathfrak{p})\leq\dim(A/\mathfrak{p}')=\dim(P)$, one necessarily has $\mathfrak{p}=\mathfrak{p}'$ \sref{0.16.1.2.2}, which completes the proof.
\end{proof}

\begin{corollary}[16.5.10]
\label{0.16.5.10}
Under the same hypotheses as in \sref{0.16.5.9}:
\begin{enumerate}
  \item[\rm{(i)}] $M_\mathfrak{p}$ is an $A_\mathfrak{p}$-module of Cohen-Macaulay.
  \item[\rm{(ii)}] We have
  \begin{equation}
  \label{0.16.5.10.1}
  \tag{16.5.10.1}
    \dim(M)=\dim_A(M/\mathfrak{p}M)+\dim_{A_\mathfrak{p}}(M_\mathfrak{p}).
  \end{equation}
\end{enumerate}
\end{corollary}

\begin{proof}
With the notations of \sref{0.16.5.9}, let $x_i'$ be the canonical images of the $x_i$ in $A_\mathfrak{p}$ $(1\leq i\leq r)$; as the $x_i'$ belong to the maximal ideal of $A_\mathfrak{p}$ and form a $M_\mathfrak{p}$-regular sequence by flatness \sref{0.15.1.14}, we have
\[
  \operatorname{prof}_{A_\mathfrak{p}}(M_\mathfrak{p})\geq r=\dim(M)\setmin\dim(M/\mathfrak{p}M)\geq\dim_{A_\mathfrak{p}}(M_\mathfrak{p})
\]
(by virtue of \sref{0.16.1.8.1}).
But taking into account \sref{0.16.4.5.1}, the three terms of this inequality are necessarily equal, whence the corollary.
\end{proof}

\oldpage[0\textsubscript{IV-1}]{39}

\begin{corollary}[16.5.11]
\label{0.16.5.11}
Let $A$ be a local Noetherian ring; suppose that there exists a Cohen-Macaulay $A$-module of finite type $M$ with support $\Spec(A)$.
Then, for every prime ideal $\mathfrak{p}$ of $A$, we have
\begin{equation}
\label{0.16.5.11.1}
\tag{16.5.11.1}
  \dim(A)=\dim(A/\mathfrak{p})+\dim(A_\mathfrak{p}).
\end{equation}
\end{corollary}

\begin{proof}
Indeed, $\Supp(M/\mathfrak{p}M)=\Spec(A/\mathfrak{p})$ $(\mathbf{0}_\mathrm{I},~1.7.5)$ and $\Supp(M_\mathfrak{p})=\Spec(A_\mathfrak{p})$; the relation \sref{0.16.5.11.1} is thus a special case of \sref{0.16.5.10.1}.
\end{proof}

\begin{corollary}[16.5.12]
\label{0.16.5.12}
Every quotient ring of a local Noetherian ring $A$ satisfying the conditions of \sref{0.16.5.11} (in particular every quotient ring of a local Cohen-Macaulay ring) is catenary.
\end{corollary}

\begin{proof}
It suffices evidently to prove that $A$ itself is catenary, that is to say for two prime ideals $\mathfrak{p}\subset\mathfrak{q}$ of $A$, we have \sref{0.16.1.4.2}
\[
  \dim(A_\mathfrak{q})=\dim(A_\mathfrak{q}/\mathfrak{p}A_\mathfrak{q})+\dim(A_\mathfrak{p}).
\]
But, by virtue of \sref{0.16.5.10}, (i) $A_\mathfrak{q}$ satisfies the same hypotheses as $A$, and the preceding relation is then none other than \sref{0.16.5.11.1} applied to $A_\mathfrak{q}$ and to the prime ideal $\mathfrak{p}A_\mathfrak{q}$ of $A_\mathfrak{q}$.
\end{proof}

\begin{env}[16.5.13]
\label{0.16.5.13}
Let $A$ be a Noetherian ring, $M$ an $A$-module of finite type.
We say that $M$ is a \emph{Cohen-Macaulay $A$-module} if, for every prime ideal $\mathfrak{p}$ of $A$, $M_\mathfrak{p}$ is an $A_\mathfrak{p}$-module of Cohen-Macaulay; by virtue of \sref{0.16.5.10} this definition coincides with \sref{0.16.5.1} when $A$ is \emph{local}.
We say that $A$ is a \emph{Cohen-Macaulay ring} if all the $A_\mathfrak{p}$ are.
It is clear that if $M$ (resp.~$A$) is a Cohen-Macaulay $A$-module (resp.~a Cohen-Macaulay ring), $S^{-1}M$ is a Cohen-Macaulay $S^{-1}A$-module (resp.~$S^{-1}A$ is a Cohen-Macaulay ring) for every multiplicative subset $S$ of $A$.
\end{env}
